\section{Introduction}

Reinforcement learning (RL)~\cite{rlbook} has become a cornerstone for advancing the reasoning capabilities of Large Language Models (LLMs), notably the Reinforcement Learning with Verifiable Reward (RLVR) paradigm~\cite{RLVR,R1}. Methods like Group Relative Policy Optimization (GRPO) \cite{GRPO, R1} have achieved remarkable success by adopting a \emph{multi-outcome} approach, generating a group of responses for each prompt to construct an on-the-fly baseline for variance reduction. While this ``group-based'' paradigm has pushed the state of the art, it suffers from fundamental inefficiencies. When all responses in a group share the same outcome (e.g., all correct or all incorrect), the relative advantage collapses to zero, yielding no learning signal. This degeneracy represents a fundamental waste of computation and data. To counteract this, a series of engineering heuristics like dynamic sampling \cite{DAPO} have been developed. These workarounds, while functional, add significant complexity and create a less principled, more convoluted optimization process.

Group-based architectural choice also imposes a critical synchronization barrier. In distributed training, the entire group must wait for its slowest member, a bottleneck that becomes particularly acute in complex agentic tasks requiring multi-turn tool use or long-horizon reasoning \cite{gao2025beyond, xu2025agents,zeng2025glm}. In these settings, interaction times are highly variable (e.g., number of interaction turns, time per interaction, etc), and a single slow-running agentic trajectory can stall its entire group, severely hindering training throughput and scalability.

We advocate for returning to the classic single-stream paradigm for policy gradient optimization~\cite{rlbook}, where each training sample is a single stream of prompt-response pair. This is not a mere simplification, but a deliberate re-alignment with foundational RL principles to address the aforementioned architectural flaws. To overcome the critical challenge of high gradient variance in this setting, we introduce Single-stream Policy Optimization (SPO). SPO replaces the noisy, on-the-fly group baseline with three synergistic components for stable and efficient learning. First, it employs a lightweight Bayesian value tracker to maintain a persistent, temporally-informed estimate of the success probability for each prompt, serving as a low-variance baseline. Second, it normalizes advantages globally across the entire batch, avoiding the instability of per-group statistics. Finally, this architecture naturally enables an adaptive curriculum via prioritized sampling, focusing computational resources on the most informative prompts.


The benefits of this principled approach are clear: SPO is inherently more scalable and eliminates the computational waste of degenerate groups. Our experiments confirm these advantages, demonstrating that SPO consistently outperforms GRPO on challenging reasoning benchmarks, improving the absolute point gains on challenging
datasets, including $7.3~\text{percentage points}~(\mathrm{pp})$ on BRUMO 25, $4.4~\mathrm{pp}$ on AIME 25, $3.3~\mathrm{pp}$ on HMMT 25, and the pass@$k$ curves of SPO are above GRPO for all $k$s. The scalability benefit is particularly pronounced in agentic settings. Our simulations, designed to model these variable-time scenarios, show that SPO's group-free design can achieve a $4.35\times$ training throughput speedup by eliminating group synchronization bottlenecks. SPO thus provides a more robust foundation for modern LLM optimization, prompting a re-evaluation of essential versus incidental complexity in the field.
